\paragraph{The polynomial on the same component.}
Let \(d=d_F\), and write \(\theta(w)\) for the smooth reference
curvature of the root combination. For a weight \(w\) and real
parameter \(s\) define \(I_R(w,s)\) by taking the arithmetic
\((d+1)\)-st power of the root line at \(Q_B\) with its original
truncated metric at \(R\), plus \(sA\), on the component \(c\).
Take the section and field degree normalizations used for
\(\Phi_B\). Since the logarithms of the determinant norms at fixed \(R\)
are linear in the weight, multilinearity gives the degree bound
\[
                      \deg_{w,s} I_R(w,s)\le d+1 .
\] The
cutoff \(R\) in the archimedean metric is independent of the
finite cutoff \(B\).

To bound this polynomial, compare its metric with the perturbation already controlled on horizontal and vertical curves. In a toroidal
chart let \(r\) again denote the logarithmic boundary radius.
The truncated matrix on the first block is the interpolation
\[
 (1-a_R)H_0+a_R H
\]
of the smooth and Hodge Gram matrices.
The same formula on the other block uses the reversed dual
matrices. For \(r\le R^2\) one has
\[
 a_R=(1+(r/R)^\kappa)^{-1}.
\]
The interpolation agrees with
the smooth reference at \(r\ge R^3\). Choose the exponent as
in the chart, \(\kappa>2\mu+2\), where each original form lies
between \(r^{-\mu}\) and \(r^\mu\) times the other. For
\(r\le\sqrt R\), the relative error in the interpolation is at
most \(R^{-1}\), since
\[
 (1-a_R)r^\mu\le R^{(\mu-\kappa)/2}.
\]

The log of each graded norm differs by \(O(1/R)\), by comparing
successive principal determinants. Everywhere its change from
the smooth reference is \(O(\log r)\). Both estimates are
uniform for bounded weights. \label{cor:flag-wall-polynomial-1}

\begin{proposition}[Componentwise arithmetic comparison at the weight wall]
\label{prop:arithmetic-wall-polynomial}
On the fixed component \(c\), the arithmetic power of the
perturbed metric is approximated by the limiting weight polynomial
from \(I_R\). Its logarithmic boundary mass satisfies
\eqref{eq:flag-wall-tail}. The difference of the finite-cutoff
powers satisfies \eqref{eq:flag-wall-finite-cutoff}. Writing \(g_c(v)\) for the geometric degree on the same component, the component measures enter as
\[
 g_c(v)=w_cG(v),\qquad w_c>0,\qquad \sum_c w_c=V.
\]

\end{proposition}

\begin{proof}
\proofstep{Comparing with the intersection in the perturbed metric.}
For the regularized potentials of Section~\ref{subsec:flag-wall-positivity}
write \(\ell=s^{-2}\). On \(r\le\ell\) their distance from the
Hodge potentials is \(O(s)\). Indeed on a ball of radius
\(c e^{-r/2}\) in a crossing chart the new boundary radius is
at most \(r+C\). Differentiating the Hodge Gram ratios bounds
the gradient there by
\[
 C(r+C)^\nu e^{r/2}.
\]

In a logarithmic
coordinate the derivative of a period brings down a factor
\(1/z_i\), bounded by \(Ce^{r/2}\), and the Gram matrices,
their inverses in the eigenframes, and the smooth changes of
frame have polynomial bounds in \(r\). The smoothing radius
\(e^{-\ell}\) times this gradient is \(O(e^{-\ell/3})\).
The bounded regularization corrections of size
\(O((1+\log\ell)/\ell)\) there are also \(O(s)\).
For \(r\ge\ell\) the upper and lower value bounds give a
difference \(O(\log r)\). \label{cor:flag-wall-polynomial-2}

\proofstep{Controlling the logarithmic boundary growth.}
Here is the boundary estimate that permits use of these comparisons
under the intersection integral. A positive curvature of the form
\(\theta+dd^c u\) off the divisor, with
\[
 -C\log r\le u\le0
\]
and \(-C\Omega\le\theta\le C\Omega\),
will be used with \(u\) smooth or tending to \(-\infty\) at the
boundary. Each curvature in the interpolation can be written
as a bounded combination of such forms. More explicitly, for
the truncated Gram metric one adds a fixed multiple of a positive
Poincar\'e form, then subtracts it.

In crossing coordinates the
latter dominates the smooth directions and the forms
\[
 \frac{i\,dz_a\wedge d\bar z_a}
 {|z_a|^2(1+|\log |z_a||)^2}
\]
in the boundary directions. Its potential is a sum of negative
logarithms of the branch radii, comparable to \(-\log r\)
up to fixed factors. The lower curvature bound for the interpolated
metric uses the blockwise decreasing weights. In the chart curvature
formula the two end metrics have bounded normalized curvature
and their connection difference is bounded in the Poincar\'e
metric. Inserting the interpolated Gram matrix in
\(\bar\partial(H^{-1}\partial H)\), the terms with a derivative
of \(a_R\) are controlled by the logarithmic derivatives of the
cutoff.

On \(r\le R^2\) these are obtained from
\[
 da_R=-\kappa a_R(1-a_R)\,dr/r.
\]
On the terminal transition the
squared smooth cutoff and \(\kappa>2\mu+2\) bound the terms
involving the inverse Gram matrix. The flag restriction adds
the nonnegative second fundamental form for each proper step.
This is the interpolation curvature estimate on the adapted
toroidal charts. The positive correction can be enlarged so its
potential dominates the absolute log growth. For the \(s\)-metric
use its semipositivity and uniform upper bound. Majorization of the
Hodge potential gives the same logarithmic lower bound off the
divisor, uniformly in \(s\). Constant changes in potentials leave
their curvatures unchanged. \label{cor:flag-wall-polynomial-3}

For a product \(d\mu\) of \(d\) of the indicated positive forms,
the mass is bounded, and for \(D\ge e\)
\begin{equation}\label{eq:flag-wall-tail}
 \int_{\{r\ge D\}}\log r\,d\mu
       \ \ll\ \frac{(1+\log D)^{d+1}}{D}.
\end{equation}
To verify the bound on a shell \(D\le r\le 2D\), truncate
each \(u\) from below by a smooth convex maximum with
\(-C\log(2D)-2\), agreeing with \(u\) on the shell.

Adding
\(C\Omega\) makes the truncated curvatures positive everywhere.
They are smooth through the divisor, since a potential tending
to \(-\infty\) is then constant near it. Their mixed masses
against powers of \(\Omega\) are uniformly bounded by Stokes'
theorem applied to the fixed smooth parts. Represent each
truncated potential up to a constant by
\(0\le v\le C'(1+\log D)\). Take a smooth increasing concave
truncation \(r_D\) of \(r\), equal to \(r\) up to \(2D\) and
constant beyond \(3D\), with derivative at most one.
The logarithmic divisor formula gives \(dd^c r_D\le C\Omega\).
Also
\[
 \int r_D\Omega^d\le\int r\Omega^d=O(1).
\]

Write one truncated curvature as \(dd^c v+\alpha\) where
\(\alpha\) is smooth and uniformly bounded above by \(C\Omega\).
For a closed mixed positive product \(P\) of complementary
degree, Stokes gives
\[
 \int r_D(dd^c v+\alpha)\wedge P
   =\int v\,dd^c r_D\wedge P+\int r_D\alpha\wedge P
   \le C'(1+\log D)\int\Omega\wedge P+
                              C\int r_D\Omega\wedge P.
\]
The first integral on the right is a uniformly bounded mass.
The second replaces one truncated curvature by \(\Omega\).
Iteration terminates at \(\int r_D\Omega^d=O(1)\), so for a mixed
product \(P_d\) of \(d\) such curvatures the shell bound is
\[
 \int_{\{D\le r<2D\}}P_d
  \ \le\ D^{-1}\int r_D P_d
  \ \ll\ (1+\log D)/D \qquad (d\ge1).
\]
Summing with weight \(\log(2D)\) on dyadic shells gives
\eqref{eq:flag-wall-tail}.
The mass on the part of bounded radius follows by the same
truncation and ordinary mixed-mass bound. This applies also
to absolute values of integrals of the bounded combinations,
by expanding their products. \label{cor:flag-wall-polynomial-4}

\proofstep{The difference of arithmetic powers.}
Changing just the archimedean metric in the arithmetic power
telescopes into \(d+1\) integrals of the potential difference
against mixed curvatures, with the common component
normalization. Write \(\varphi_{R'},\varphi_R\) for the normalized root
potentials in the two readings and \(\mathcal M_i\) for the products
of \(i\) curvatures from \(R'\) and \(d-i\) from \(R\).
Use the \(O(1/R)\) bound on \(r\le\sqrt R\) and the weighted
tail \eqref{eq:flag-wall-tail} on its complement:
\[
 \sum_{i=0}^d\left|\int (\varphi_{R'}-\varphi_R)\,\mathcal M_i\right|
   \ll R^{-1}+\frac{(1+\log R)^{d+1}}{\sqrt R}.
\]
The same component normalization is used in the two degrees,
and the bound is uniform in \(R'\ge R\), \(B\), the bounded
dominant weights and \(0\le s\le1\). Thus they converge on an
open set of strictly decreasing weights and \(0<s<1\).
Evaluation on that open set separates polynomials of bounded
degree.

A finite subset of evaluations forms a basis of their
dual, so that all the coefficients converge. Call the limiting
polynomial \(P_{B,c}\). At \(s=0\) its component sum is
\(\Phi_B(L_w^{d+1})\) by the defining limit, first on the open
chamber and then identically. Compare the \(s\)-metric to the
truncated metric at \(R\ge\ell^2\). Their difference is \(O(s)\)
on \(r\le\ell\) and \(O(\log r)\) elsewhere. Putting \(\ell=s^{-2}\) in the tail estimate gives
\[
 \ell^{-1}(1+\log\ell)^{d+1}
 =s^2(1+2\log(1/s))^{d+1}=O(s)
\]
as \(s\downarrow0\), for the fixed dimension \(d\). Taking \(R\to\infty\)
proves that the actual component intersection for \(M_B(v,s)\)
differs from
\[
 P_{B,c}(v-sd_0,s)
\]
by \(O(s)\), uniform in
\(B\) on the fixed model. \label{cor:flag-wall-polynomial-5}

\paragraph{Dependence on the finite cutoff.}

For wall estimates obtained by varying the cutoff to determine the functional at a fixed cutoff, the difference of the two arithmetic degrees must be tracked prime by prime.

Compare \(B\) with \(2\) on a common refinement and divide by
the corresponding generic cover degrees. The changed primes
are good and their root log norms differ by at most
\(\log p/(p-1)\) in the normalization at \(p\). For a combination
\(x\) at \(B\) and \(y\) at \(2\) with the same bounded weights
and ample summand, the generic and archimedean lines coincide.
The vertical log bound for \(x-y\) is at most the sum of the
absolute weight coefficients times the root bound. On a list
of \(d\) lines nef over the changed primes its mixed intersection
is therefore bounded in absolute value by
\[
 C\sum_{p\le B}\log p/(p-1)
\]
times their generic mixed degree.
Indeed the local vertical divisor is between the two corresponding
multiples of the whole fibre. Restrict the nef lines to the
components and sum their nonnegative mixed degrees with these
multiplicities.

For clarity, \(y\) itself need not be nef at the newly changed
primes. On the model at \(2\) take a fixed relatively ample
\(H\) large enough that \(H\pm x_i\) are relatively nef for
every reference root \(x_i\). For \(\|w\|_\infty\le1\) and
\(s\ge0\), \(y+h\) is in the nef cone, where \(h\) is the
number of roots times \(H\). This follows by expanding each
\(H-w_i x_i\) with coefficients \((1\pm w_i)/2\).

At the changed primes \(x\) is nef by the Hodge--Tate step
comparison. Use
\[
 x^{d+1}-y^{d+1}=\sum_{i=0}^d(x-y)x^i y^{d-i},\qquad
 y^{d-i}=\sum_{j=0}^{d-i}\binom{d-i}{j}(-1)^j
                                 (y+h)^{d-i-j}h^j .
\]
After taking absolute values of these coefficients, every
local mixed degree uses only the nef lines \(x,y+h,h\).
In the bound by generic degrees replace \(x\) by \(y\),
whose generic class agrees. The sum over \(j\) is
\(\deg_{\mathrm{geom}}(y^i(y+2h)^{d-i})\). It is bounded on the
fixed compact coefficient range independently of the refinement,
after division by the cover degree. This yields
\begin{equation}\label{eq:flag-wall-finite-cutoff}
 \lvert I_R(w,s;B)-I_R(w,s;2)\rvert\ll\log B .
\end{equation}
The same inequality holds in the limit. Applying the same finite matrix of evaluations on the open
chamber now bounds every coefficient of \(P_{B,c}\) by \(O(\log B)\).
\label{cor:flag-wall-polynomial-6}

\paragraph{The geometric degree.}
The canonical root class \(L_v\) is nef for dominant \(v\).
Indeed the regularizations above approximate it by smooth forms
with curvature at least \(-\varepsilon\Omega\), with
\(\varepsilon\downarrow0\). The shifted geometric class
\(M=L_{v-sd_0}+s A\) is nef for dominant \(v-sd_0\). In the
polarization of the toroidal model \(L_{d_0}-A\) has the class
of an effective boundary divisor \(D_0\), with the multiplicities
of its local boundary section. The telescoping formula gives
\[
  \deg L_v^d-\deg M^d
     =s\sum_{i=0}^{d-1} D_0\cdot L_v^i M^{d-1-i}\ge0 .
\]
Each summand is nonnegative by restriction to the effective
components and approximation by positive forms. Polynomial
continuity gives convergence to \(\deg L_v^d\) as \(s\downarrow0\).

For the canonical extension of the homogeneous bundles over a
smooth toroidal compactification of a neat arithmetic quotient,
Mumford's proportionality theorem identifies the Chern evaluation
with the compact-dual evaluation times the component volume in
the proportionality normalization. Passing to the flags uses the
flag-bundle pushforward on both sides, so the weighted root
powers have the same factor. Dividing by the degree of a common
neat cover gives the same formula on the selected component.
It holds first at integral weights and then identically by
multilinearity. The resulting upper degree is \(w_cG(v)\),
where \(w_c>0\) and \(\sum_c w_c=V\).
\cite{mumford-flag-wall}\label{cor:flag-wall-polynomial-7}
\end{proof}
